\documentclass{beamer}

\input{../preambulo/inputenc.tex}
\input{../preambulo/fuente.tex}
\input{../preambulo/spanish.tex}
\input{../preambulo/tikz.tex}
\input{../preambulo/matematica.tex}
\input{../preambulo/colors.tex}
\input{../preambulo/flotantes.tex}
\input{../preambulo/beamer_zub.tex}
\input{../preambulo/paquetes.tex}
\usepackage{ccicons}
% \ccPublicDomain

\newcommand{\etal}{{\it et al.}}
\newcommand{\ie}{{\it i. e.}}
\newcommand{\zlink}[1]{\href{#1}{\color{celeste}\faExternalLink*}}


\graphicspath{{../img/}{img/}}
\framebg{bg/pizarron3.png} 
\hypersetup{colorlinks=true, urlcolor=azul!40!white}

\begin{document}

\newcommand\CC{}

\begin{zframe}{}% yscale=0.9, % with frametitle
\path(cp) ++(3,0) node(I)[]{\includegraphics[angle=90,width=6cm]{img/35.png}};
\path(c2|-I.north) ++(0,-.5) node(A)[anchor=north,align=left]{
  \color{verde} \large\textbf{Métodos Computacionales}\\[3mm]  
  \color{celeste} \textbf{Modelado Molecular}\\[2mm]  
  \color{lila} \textit{Octubre 22, 2025}
};
\normalsize
\path(A.west|-I.south) ++(0,.5) node(B)[anchor=south west,align=left]{
  S. Alexis Paz\\[5mm]

\includegraphics[width=3cm]{logos/DQTC_orange.png}};

\path(I.south east) ++(0,0) node[anchor=north east,inner sep=0]{
  \tiny Bernal \textit{et al.}, \href{https://doi.org/10.1098/rspa.1964.0147}{Proc. Roy. Soc. A 280(1964)16}};
 
\end{zframe}

\renewcommand\CC{
  \path(se) node[anchor=south east]{\tiny\color{gray} MC2025 - S.A.Paz};}

\begin{zframe}{}

\path(t-|np) ++(-.5,.1) node(last)[def,anchor=north,align=left,text width=10cm]{\\
Sistema lógico y matemático, construido a partir de un conjunto de postulados y
conceptos fundamentales, cuyo propósito es explicar y predecir el comportamiento de un
sistema físico dentro de un determinado dominio de aplicación.};
\path(last.north west) node(H)[defhead,anchor=west]{Teoría}; 
                                                        
\node[zf/text=next,shift={(0,-1)},def,text width=10cm]{\\[1mm]
  Representación simplificada y controlada de un sistema
  físico real, construida con el propósito de entender, describir o predecir su
  comportamiento, ya sea mediante experimentos, análisis matemático o
  simulación.};
\path(last.north west) node(H)[defhead,anchor=west]{Modelo}; 
                                                         
% \node[zf/text=next,def,text width=10cm]{\\[1mm]
%   Representación matemática y computacional del sistema físico que se desea
%   estudiar, construida con el objetivo de describir su comportamiento a nivel
%   atómico o molecular};
% \path(last.north west) node(H)[defhead,anchor=west]{Simulación}; 
                
\path(ne) ++(.2,-.1) node(last)[anchor=north east,rotate=-10]{\includegraphics[width=2cm]{bib/Cramers.png}};
\path(last.south) ++(0,.1) node[anchor=north,rotate=10]{\includegraphics[width=2cm]{bib/Allen.png}};
%
% \node[zf/text=next, text width=12cm]{
% Representación matemática y computacional del sistema físico
% que se desea estudiar, construida con el objetivo de describir su
% comportamiento a nivel atómico o molecular.};

\end{zframe}


\begin{zframe}{yscale=0.9}
\frametitle{Modelos moleculares}

\path<1>(cp) ++(0,0) node(A){\includegraphics[width=10cm]{33.png}};
\path<2>(cp) ++(0,0) node(A){\includegraphics[width=10cm]{34.png}};
\path<3>(cp) ++(0,0) node(A){\includegraphics[width=10cm]{32.png}};
\path(A.south west) ++(0,0) node[anchor=north west,inner sep=0]{
  \tiny Bernal \textit{et al.}, \href{https://doi.org/10.1098/rspa.1964.0147}{Proc. Roy. Soc. A 280(1964)16}};
\end{zframe}
  
\begin{zframe}{}
                          
\path<1>(cp) ++(0,0) node(A){\includegraphics[width=10cm]{36.png}};
\path<2>(cp) ++(0,0) node(A){\includegraphics[width=10cm]{37.png}};
\path(A.south west) ++(0,0) node[anchor=north west,inner sep=0]{
  \tiny Bernal \textit{et al.}, \href{https://doi.org/10.1098/rspa.1964.0147}{Proc. Roy. Soc. A 280(1964)16}};
                    
\end{zframe}
             
\begin{zframe}{
}
                          
\path(hl) ++(.5,0) node(A){\includegraphics[width=7cm]{29b.png}};
\path(A.south west) ++(.1,0) node[anchor=north west,inner sep=0]{
  \tiny Bernal \textit{et al.}, \href{https://doi.org/10.1098/rspa.1964.0147}{Proc. Roy. Soc. A 280(1964)16}};
\path(hr) ++(.5,0) node(A){\includegraphics[width=5cm]{31.png}};
\path(A.south west) ++(.1,0) node[anchor=north west,inner sep=0]{
  \tiny Bernal \textit{et al.}, \href{https://doi.org/10.1098/rspa.1964.0147}{Proc. Roy. Soc. A 280(1964)16}};
                    
\end{zframe}
 
\begin{zframe}{yscale=0.9}
\frametitle{Mecánica Molecular}
 
\node[zf/text=first,shift={(0,.4)}]{\Large Mecánica Clásica $\rightarrow$ \color{amarillo} Born-Oppenheimer};
\draw[verde,ultra thick,<-](last.west) -- ++(-.5,0) -- ++(0,.5); 

\node[zf/text=next]{
${\ds \hat H=\color{celeste}-\sum_{i}\frac {\nabla _{i}^{2}}{2}-\sum_{i,A}\frac {Z_{A}}{r_{iA}}+\sum_{i>j}\frac {1}{r_{ij}}\color{naranja}-\sum_{A}\frac{\nabla _{A}^{2}}{2M_{A}}+\sum_{B>A}\frac {Z_{A}Z_{B}}{R_{AB}}}$\\[2mm]
${\ds \hat H=\color{celeste}\hat H_e(R,r)\color{naranja}-\sum _{A}{{\frac {\nabla _{A}^{2}}{2M_{A}}}}+\sum _{B>A}{\frac {Z_{A}Z_{B}}{R_{AB}}}}\quad{\ds \Psi(R,r)=\color{celeste}\psi(R,r)\color{naranja}\phi(R)}$
};
    
\node[zf/text=next,shift={(0,.3)},text width=11cm]{
\begin{align*}
\ds \text{Schrödinger: }\quad \hat H\Psi(R,r)&=E\Psi(R,r)\\
\ds {\color{celeste}\psi(R,r)}\ab[{\color{celeste}E_e(R)}{\color{naranja}-\sum _{A}\frac{\nabla_{A}^{2}}{2M_{A}}+\sum_{B>A}\frac {Z_{A}Z_{B}}{R_{AB}}}]{\color{naranja}\phi(R)}&={\color{celeste}\psi(R,r)}E{\color{naranja}\phi(R)}\\
\ds \ab[ {\color{naranja}-\sum_{A}\frac{\nabla_{A}^{2}}{2M_{A}}}+{\color{verde}U(R)} ]{\color{naranja}\phi(R)}&=E\color{naranja}\phi(R)\\
\end{align*}
};
       
                   
\end{zframe} 
                  
\begin{zframe}{yscale=0.9}
\frametitle{Mecánica Molecular}
                          
\path(c1) ++(0,0) node(A){\includegraphics[width=4cm]{CAS1.png}};
\path(c4) ++(0,0) node(A){\includegraphics[width=4cm]{DOS.png}};

\path(hl) ++(.5,0) node(A){\includegraphics[width=6cm]{DSC02426.JPG}};
                    
\end{zframe}
              
\begin{zframe}{} %yscale=0.9

\path(t) ++(0,0) node(T)[align=left,anchor=north west]{
{\color{verde} \Large Potencial de a pares:}
};
                                

\path(t|-T.south) ++(0,-.1) node(Tx)[align=left,anchor=north west]{
${\ds U(\vb{r})=\sum_{i}U_{1}(\vb{x_i})+\sum_{i<j}U_{2}(\vb{x_i},\vb{x_j})+\sum_{i<j<k}U_{3}(\vb{x_i},\vb{x_j},\vb{x_k})+\dotsb \label{e_exppot}}$\\[2mm]
${\ds U(\vb{r})=\sum_{i<j}U_{2}(r_{ij})+\sum_{i<j<k}U_{3}(\vb{r}_{ij},\vb{r}_{jk})+\dotsb \label{e_exppot}}$\\[2mm]
};
    
\path(t|-Tx.south) ++(0,-.2) node(T)[align=left,anchor=north west]{
\color{celeste}Lennard-Jones\\[2mm] %\cite{Jones1924,Jones1931} %\cite{Rahman1964,Alder1959}
${\ds U_2(r_{ij})=4\epsilon\ab[\ab(\frac{\sigma}{r_{ij}})^{12}-\ab(\frac{\sigma}{r_{ij}})^{6}]}$
};
       
     
\path(t|-T.south) ++(0,-.2) node(T)[align=left,anchor=north west]{
\color{celeste}Morse\\[2mm]
${\ds U_2(r_{ij})=D_e(1-e^{\alpha(r_{ij}-r_e)})^2-D_e}$
};
                               
\path(cp|-Tx.south) ++(2,-.2) node(T)[align=left,anchor=north west]{
\color{celeste}Fuerza\\[2mm]
${\ds \vb{F}_{r_i}=\left.\nabla_{r_i} U_2(r_{ij})\right|_{r1}}$\\[2mm]
${\ds \quad=\pdv{U_2(r_{ij})}{r_{ij}}\nabla_{r_i}r_{ij}}$\\[2mm]
${\ds \quad=\pdv{U_2(r_{ij})}{r_{ij}}\vu{r}_{ij}}$
};
                     
\end{zframe} 
         
\begin{zframe}{} %yscale=0.9

\path(t) ++(0,0) node(T)[align=left,anchor=north west]{
{\color{verde} \Large Campo de Fuerzas:}
};
                                
\path(t|-T.south) ++(-.2,-.2) node(T)[align=left,anchor=north west]{
\color{celeste}Hill (1942)\\[2mm]
${\ds U\ab(\ab\{r_{ij}\}, \ab\{d\}, \ab\{\phi\})=U^{\text{LJ}}_2(r_{ij}) + \sum_{d} k_d\ab(d-d_{eq})^2 + \sum_{\phi} k_\phi\ab(\phi-\phi_{eq})^2}$
};

\path(t|-T.south) ++(-.2,-.2) node(T)[align=left,anchor=north west]{
\color{celeste}OPLS (1984)\\[2mm]
${\ds U\ab(\ab\{r_{ij}\},\ab\{d\},\ab\{\phi\},\ab\{\theta\})=\sum_{i<j}f_{ij}\ab(\frac{q_iq_je_0^2}{r_{ij}}+U^{\text{LJ}}_2(r_{ij}))}$ \\[2mm]
${\ds +\sum_{d} k_d\left(d-d_{eq}\right)^2 + \sum_{\phi} k_\phi\left(\phi-\phi_{eq}\right)^2 + \sum_{\theta} \sum_{j=0}^{5} C_j\cos^j(\theta)}$
};
   
\path(t|-T.south) ++(-.2,-.2) node(T)[align=left,anchor=north west]{
\color{celeste}CHARMM, GROMOS, AMBER, Tersoff, ReaxFF, \ldots
};

\path(ep) ++(1.2,-.5) node(A)[anchor=east]{\includegraphics[angle=150,width=5.6cm]{dihedrals.png}};
                   
\end{zframe} 
           
\begin{zframe}{} %yscale=0.9

\path(t) ++(0,0) node(T)[align=left,anchor=north west]{
{\color{verde} \Large Muchos Cuerpos:}
};
                                
\path(t|-T.south) ++(-.2,-.2) node(T)[align=left,anchor=north west]{
\color{celeste}El método del átomo embebido, EAM (1984)\\[2mm]
${\ds  U(\vb{r})=\frac{1}{2}\sum_{i<j}U_2(r_{ij})+\sum_{i}E_i\ab(\rho_i\ab(\ab\{r_{ij}\}))}$
};

\path(t|-T.south) ++(1,-.2) node(T)[align=left,anchor=north west]{
${\rho_i=\sum_{j\neq i}\psi_j\ab(r_{ij})}$
};
                   
           
\path(vl) ++(0,0) node(A){\includegraphics[width=9cm]{SS_rot.png}};
\path(c1) ++(.3,-.7) node(A){\includegraphics[width=5cm]{F4_PCCP.png}};
                    
\end{zframe}
            
\begin{zframe}{}
         
\path<-2>(t) ++(0,.5) node(T)[align=left,anchor=north west]{
{\color{verde} \Large Lista de Vecinos}
};                                    
                           
\path<-2>(t-|c1) ++(0,.5) node(T)[align=left,anchor=north west]{
\color{celeste} Escalea con $N^{2}$
};                                    
          
\path<3->(t) ++(0,.5) node(T)[align=left,anchor=north west]{
{\color{verde} \Large Lista de Vecinos - Verlet}
};                                    
                          
\path<3->(t-|c1) ++(0,.5) node(T)[align=left,anchor=north west]{
\color{celeste} Escalea con $N^{5/3}$
};                                    
                  
\path<4->(2.2,1) node[align=left]{Actualizo si por ejemplo\\$\max\ab(\ab|r_n-r_0|)>d/2$};

\path(5,2) coordinate(scope);
\begin{scope}[x=1cm,y=1cm,amarillo,shift=(scope),thick]

% % Definir
% \path(0,0) coordinate(O);
% \path(5,5) coordinate(K1);          % Esquina del cuadrante 1
% \path(K1-|{(0,5)}) coordinate(K2); % Esquina del cuadrante 2
% \path(K2|-{(0,0)}) coordinate(K3); % Esquina del cuadrante 3
%
% % Emergentes
% \path(K3-|K1) coordinate(K4);       % Esquina del cuadrante 4
% \path(O-|K1) coordinate(xp);       % Punta de eje +x
% \path(O-|K2) coordinate(xn);       % Punta de eje -x
% \path(O|-K1) coordinate(yp);       % Punta de eje +y
% \path(O|-K3) coordinate(yn);       % Punta de eje -y


% grid
% \draw[ystep=5, xstep=5, celeste] (K3) grid (K1);

\fill(-3,3)  circle(.1) coordinate(V);
\fill(2,2)  circle(.1) coordinate(W);
\fill(1,.2) circle(.1) coordinate(X);
\fill(3,4)  circle(.1) coordinate(Y);
\fill(3.6,.7)  circle(.1) coordinate(Z);
     
\draw[dashed](V) circle(1);
\draw[dashed](W) circle(1);
\draw[dashed](X) circle(1);
\draw[dashed](Y) circle(1);
\draw[dashed](Z) circle(1);
\draw[thick](V) -- ++(20:1) node[pos=.5,above]{$\sigma$};
\draw<2->[rojo,thick](V) -- ++(120:2.5) node[pos=.5,above right]{$r_{\text{cut}}$};
\draw<3->[lila,thick](V) -- ++(-90:3) node[pos=.5,right]{$d$};
             
\draw<2->[rojo,dashed](V) circle(2.5);
\draw<2->[rojo,dashed](W) circle(2.5);
\draw<2->[rojo,dashed](X) circle(2.5);
\draw<2->[rojo,dashed](Y) circle(2.5);
\draw<2->[rojo,dashed](Z) circle(2.5);
                 
\draw<3->[lila,dashed](V) circle(3);
\draw<3->[lila,dashed](W) circle(3);
\draw<3->[lila,dashed](X) circle(3);
\draw<3->[lila,dashed](Y) circle(3);
\draw<3->[lila,dashed](Z) circle(3);

\draw<4->[verde](Z) -- (W);
\draw<4->[verde](Z) -- (X);
\draw<4->[verde](X) -- (W);
\draw<4->[verde](W) -- (Y);
                 
\end{scope} 
      
\end{zframe}
                        
\begin{zframe}{}
         
\path(t) ++(0,.5) node(T)[align=left,anchor=north west]{
{\color{verde} \Large Celdas Vecinas}
};                                    
                         
\path(t-|c1) ++(0,.5) node(T)[align=left,anchor=north west]{
\color{celeste} Escalea con $N$
};                                    
                  
\path(3,2) coordinate(scope);
\begin{scope}[x=1cm,y=1cm,amarillo,shift=(scope),thick]

% Definir
\path(-2,-2) coordinate(O);
\path(9,5.8) coordinate(K1);          % Esquina del cuadrante 1
\path(K1-|{(-1.8,5.8)}) coordinate(K2); % Esquina del cuadrante 2
\path(K2|-{(-1.8,-1.8)}) coordinate(K3); % Esquina del cuadrante 3
 
% Emergentes
\path(K3-|K1) coordinate(K4);       % Esquina del cuadrante 4
\path(O-|K1) coordinate(xp);       % Punta de eje +x
\path(O-|K2) coordinate(xn);       % Punta de eje -x
\path(O|-K1) coordinate(yp);       % Punta de eje +y
\path(O|-K3) coordinate(yn);       % Punta de eje -y


% grid
\draw[ystep=2.5, xstep=2.5, celeste] (K3) grid (K1);

\fill(-1,3)  circle(.1) coordinate(V);
\fill(4,2)  circle(.1) coordinate(W);
\fill(3,.2) circle(.1) coordinate(X);
\fill(5,4)  circle(.1) coordinate(Y);
\fill(5.6,.7)  circle(.1) coordinate(Z);
     
\draw[dashed](V) circle(1);
\draw[dashed](W) circle(1);
\draw[dashed](X) circle(1);
\draw[dashed](Y) circle(1);
\draw[dashed](Z) circle(1);
\draw[thick](V) -- ++(20:1) node[pos=.5,above]{$\sigma$};
\draw[rojo,thick](V) -- ++(120:2.5) node[pos=.5,above right]{$r_{\text{cut}}$};
\draw[<->,ultra thick,celeste](5,3) -- (7.5,3) node[pos=.5,below]{$r_{\text{cut}}$};
         
\draw[rojo,dashed](V) circle(2.5);
\draw[rojo,dashed](W) circle(2.5);
\draw[rojo,dashed](X) circle(2.5);
\draw[rojo,dashed](Y) circle(2.5);
\draw[rojo,dashed](Z) circle(2.5);
                
\end{scope} 
      
\end{zframe}
      
            
\begin{zframe}{}
         
\path(t) ++(0,0) node(T)[align=left,anchor=north west]{
{\color{verde} \Large Configuración Inicial:}
};                                    


\path<-3>(c2) ++(-1.5,-.5) node(A)[naranja]{$U(\vb{r})$};
\path<3->(3,1.4) node(A)[naranja]{$\approx U(\vb{r})$};


\draw<3,4>[ultra thick,naranja](1,4) -- (3,2) -- (5,4) -- (3,5) -- cycle;
\fill<3,4>[azul,yshift=4](2,3) -- (3,2) -- (4,3) -- (3,3.75) -- cycle;
\fill<3,4>[lila,yshift=7](4,3) -- (5,4) -- (4,4.5) -- (3,3.75)  -- cycle;
\fill<3,4>[amarillo,yshift=15](4,4.5) -- (3,3.75) -- (2,4.5) -- (3,5) -- cycle;
\fill<3,4>[naranja,yshift=25](3,3.75) -- (2,4.5) -- (1,4) -- (2,3) -- cycle;
 
\path<1-3>(hl) ++(0,0) node(A){\includegraphics[width=4cm]{CAS1.png}};
\draw<1>[<-,ultra thick,verde](3.8,6.2) -- ++(90:1);
\draw<2->[<-,ultra thick,verde](2.5,3.3) -- ++(-120:1);
 



\path(6,2) coordinate(scope);
\begin{scope}[x=1cm,y=1cm,amarillo,shift=(scope),thick]

% Definir
\path(0,0) coordinate(O);
\path(5,5) coordinate(K1);          % Esquina del cuadrante 1
\path(K1-|{(0,5)}) coordinate(K2); % Esquina del cuadrante 2
\path(K2|-{(0,0)}) coordinate(K3); % Esquina del cuadrante 3

% Emergentes
\path(K3-|K1) coordinate(K4);       % Esquina del cuadrante 4
\path(O-|K1) coordinate(xp);       % Punta de eje +x
\path(O-|K2) coordinate(xn);       % Punta de eje -x
\path(O|-K1) coordinate(yp);       % Punta de eje +y
\path(O|-K3) coordinate(yn);       % Punta de eje -y


% grid
\draw[ystep=5, xstep=5, celeste] (K3) grid (K1);
\fill(2,2)  circle(.1) coordinate(W);
\fill(1,.2) circle(.1) coordinate(X);
\fill(3,4)  circle(.1) coordinate(Y);
\fill<1>(2.5,.7)  circle(.1) coordinate(Z);
\fill<2->(3.6,.7)  circle(.1) coordinate(Z);
      
\draw<1>[ultra thick,dashed,naranja](W) circle(1);
\draw<1>[ultra thick,dashed,naranja](X) circle(1);
\draw<1>[dashed](Y) circle(1);
\draw<1>[ultra thick,dashed,naranja](Z) circle(1);
      
\draw<2->[dashed](W) circle(1);
\draw<2->[dashed](X) circle(1);
\draw<2->[dashed](Y) circle(1);
\draw<2->[dashed](Z) circle(1);

\draw[thick](Y) -- ++(20:1) node[pos=.5,above]{$r_{\text{cut}}$};
 
\end{scope} 
      

\end{zframe}
 
\begin{zframe}{}
         
\path(t-|cp) ++(0,0) node(A){\includegraphics[width=9cm]{lmartinez/packmol.png}};

\path(c3) ++(0,-.4) node(A){\includegraphics[width=5.5cm]{lmartinez/bilayer2.jpg}};
\path(cp) ++(0,.5) node(A){\includegraphics[width=10cm]{lmartinez/mfig002.jpg}};
\path(c4) ++(0,-.4) node(A){\includegraphics[width=5.5cm]{lmartinez/mixture2.jpg}};

\path(sw) ++(1,.1) node[anchor=south west,inner sep=2]{
  \scriptsize Martinez  \textit{et al.}, \href{https://doi.org/10.1002/jcc.21224}{J. Comp. Chem. 30(2009)2157}};
\path(sw) ++(1,.4) node[anchor=south west,inner sep=2]{
  \scriptsize Martinez \textit{et al.}, \href{https://doi.org/10.1002/jcc.10216}{J. Comp. Chem. 24(2003)819}};
 
                    
\end{zframe}
                

\begin{zframe}{}
          
\path(t) ++(0,.5) node(T)[align=left,anchor=north west]{
{\color{verde} \Large Condiciones de Contorno \color{naranja} Periódicas}
};                                    
  
\path(3,2) coordinate(scope);
\begin{scope}[x=1cm,y=1cm,amarillo,shift=(scope), thick]

% Definir
\path<1>(0,0) coordinate(O);
\path<1>(4,4) coordinate(K1);          % Esquina del cuadrante 1
\path<1>(K1-|{(0,4)}) coordinate(K2); % Esquina del cuadrante 2
\path<1>(K2|-{(0,0)}) coordinate(K3); % Esquina del cuadrante 3
\path<2>(-2,-2) coordinate(O);
\path<2>(9,5.8) coordinate(K1);          % Esquina del cuadrante 1
\path<2>(K1-|{(-1.8,5.8)}) coordinate(K2); % Esquina del cuadrante 2
\path<2>(K2|-{(-1.8,-1.8)}) coordinate(K3); % Esquina del cuadrante 3


% Emergentes
\path(K3-|K1) coordinate(K4);       % Esquina del cuadrante 4
\path(O-|K1) coordinate(xp);       % Punta de eje +x
\path(O-|K2) coordinate(xn);       % Punta de eje -x
\path(O|-K1) coordinate(yp);       % Punta de eje +y
\path(O|-K3) coordinate(yn);       % Punta de eje -y


% grid
\draw[ystep=4, xstep=4, ultra thick, celeste] (K3) grid (K1);

\fill(2,2)    circle(.1) coordinate(W);
\fill(1,.2)   circle(.1) coordinate(X);
\fill(3,3.8)  circle(.1) coordinate(Y);
\fill(3.6,.7) circle(.1) coordinate(Z);
\fill<2>(W) ++(0,-4) circle(.1) coordinate(W2);
\fill<2>(Y) ++(0,-4) circle(.1) coordinate(Y2);
\fill<2>(X) ++(0,4) circle(.1) coordinate(X2);
\fill<2>(Z) ++(0,4) circle(.1) coordinate(Z2);
\draw<2>[dashed](W2) circle(1);
\draw<2>[dashed](X2) circle(1);
\draw<2>[dashed](Y2) circle(1);
\draw<2>[dashed](Z2) circle(1);
\draw[dashed](W) circle(1);
\draw[dashed](X) circle(1);
\draw[dashed](Y) circle(1);
\draw[dashed](Z) circle(1);
\draw[thick](Y) -- ++(20:1) node[pos=.5,above]{$r_{\text{cut}}$};

\fill<2>(W) ++(4,0) circle(.1) coordinate(W);
\fill<2>(X) ++(4,0) circle(.1) coordinate(X);
\fill<2>(Y) ++(4,0) circle(.1) coordinate(Y);
\fill<2>(Z) ++(4,0) circle(.1) coordinate(Z);
\fill<2>(W) ++(0,-4) circle(.1) coordinate(W2);
\fill<2>(Y) ++(0,-4) circle(.1) coordinate(Y2);
\fill<2>(X) ++(0,4) circle(.1) coordinate(X2);
\fill<2>(Z) ++(0,4) circle(.1) coordinate(Z2);
\draw<2>[dashed](W2) circle(1);
\draw<2>[dashed](X2) circle(1);
\draw<2>[dashed](Y2) circle(1);
\draw<2>[dashed](Z2) circle(1);
\draw<2>[dashed](W) circle(1);
\draw<2>[dashed](X) circle(1);
\draw<2>[dashed](Y) circle(1);
\draw<2>[dashed](Z) circle(1);

\fill<2>(W) ++(4,0) circle(.1) coordinate(W);
\fill<2>(X) ++(4,0) circle(.1) coordinate(X);
\fill<2>(Y) ++(4,0) circle(.1) coordinate(Y);
\fill<2>(Z) ++(4,0) circle(.1) coordinate(Z);
\fill<2>(W) ++(0,-4) circle(.1) coordinate(W2);
\fill<2>(Y) ++(0,-4) circle(.1) coordinate(Y2);
\fill<2>(X) ++(0,4) circle(.1) coordinate(X2);
\fill<2>(Z) ++(0,4) circle(.1) coordinate(Z2);
\draw<2>[dashed](W2) circle(1);
\draw<2>[dashed](X2) circle(1);
\draw<2>[dashed](Y2) circle(1);
\draw<2>[dashed](Z2) circle(1);
\draw<2>[dashed](W) circle(1);
\draw<2>[dashed](X) circle(1);
\draw<2>[dashed](Y) circle(1);
\draw<2>[dashed](Z) circle(1);

\fill<2>(W) ++(-12,0) circle(.1) coordinate(W);
\fill<2>(X) ++(-12,0) circle(.1) coordinate(X);
\fill<2>(Y) ++(-12,0) circle(.1) coordinate(Y);
\fill<2>(Z) ++(-12,0) circle(.1) coordinate(Z);
\fill<2>(W) ++(0,-4) circle(.1) coordinate(W2);
\fill<2>(Y) ++(0,-4) circle(.1) coordinate(Y2);
\fill<2>(X) ++(0,4) circle(.1) coordinate(X2);
\fill<2>(Z) ++(0,4) circle(.1) coordinate(Z2);
\draw<2>[dashed](W2) circle(1);
\draw<2>[dashed](X2) circle(1);
\draw<2>[dashed](Y2) circle(1);
\draw<2>[dashed](Z2) circle(1);
\draw<2>[dashed](W) circle(1);
\draw<2>[dashed](X) circle(1);
\draw<2>[dashed](Y) circle(1);
\draw<2>[dashed](Z) circle(1);
           
\end{scope} 
      

\end{zframe}
 
\begin{zframe}{}
          
\path<1-5>(t) ++(0,.5) node(T)[align=left,anchor=north west]{
{\color{verde} \Large Condiciones de Contorno \color{naranja} Periódicas}
};                                    
                                  
\path<6>(t) ++(0,.5) node(T)[align=left,anchor=north west]{
{\color{verde} \Large Conveción de Imagen Mínima}
};                                    
                                   
\path<7>(t) ++(0,.5) node(T)[align=left,anchor=north west]{
{\color{verde} \Large Partículas Fantasma}
};                                    
                                   
\path<7>(vu) ++(1,.8) node(T)[align=left,anchor=north west]{
\color{celeste} ¡Paralelización!
};                                    
                                           
\path(3,2) coordinate(scope);
\begin{scope}[x=1cm,y=1cm,amarillo,shift=(scope),thick]

\clip<7>(-1,-1) rectangle (5,5);
\draw<7>[<->,ultra thick,celeste](4,3) -- (5,3) node[pos=.5,below]{$r_{\text{cut}}$};

% Definir
\path(-2,-2) coordinate(O);
\path(9,5.8) coordinate(K1);          % Esquina del cuadrante 1
\path(K1-|{(-1.8,5.8)}) coordinate(K2); % Esquina del cuadrante 2
\path(K2|-{(-1.8,-1.8)}) coordinate(K3); % Esquina del cuadrante 3


% Emergentes
\path(K3-|K1) coordinate(K4);       % Esquina del cuadrante 4
\path(O-|K1) coordinate(xp);       % Punta de eje +x
\path(O-|K2) coordinate(xn);       % Punta de eje -x
\path(O|-K1) coordinate(yp);       % Punta de eje +y
\path(O|-K3) coordinate(yn);       % Punta de eje -y


% grid
\draw[ystep=4, xstep=4, ultra thick, celeste] (K3) grid (K1);

\fill(2,2)    circle(.1) coordinate(W);
\fill[rojo](1,.2)   circle(.1) coordinate(X);
\fill[lila](3,3.8)  circle(.1) coordinate(Y);
\fill[naranja](3.6,.7) circle(.1) coordinate(Z);
\fill(W) ++(0,-4) circle(.1) coordinate(W2);
\fill[lila](Y) ++(0,-4) circle(.1) coordinate(Y2);
\fill(X) ++(0,4) circle(.1) coordinate(X2);
\fill[naranja](Z) ++(0,4) circle(.1) coordinate(Z2);
\draw[dashed](W2) circle(1);
\draw[dashed](X2) circle(1);
\draw[ultra thick, lila,dashed](Y2) circle(1);
\draw[ultra thick, naranja,dashed](Z2) circle(1);
\draw[dashed](W) circle(1);
\draw[ultra thick, rojo](X) circle(1);
\draw[ultra thick, lila](Y) circle(1);
\draw[ultra thick, naranja](Z) circle(1);
% \draw[thick](Y) -- ++(20:1) node[pos=.5,above]{$\sigma$};

\draw<2-6>[verde,->,ultra thick,dashed](X) -- (Z);
\draw<3-6>[verde,<-,ultra thick](Z) ++(-4,0) -- (X);
\draw<3-6>[celeste,<-,ultra thick,dashed](Z) ++(-4,0) -- (Z);
        
\draw<4-6>[verde,->,ultra thick,dashed](Z) -- (Y);
\draw<5-6>[verde,<-,ultra thick](Y) ++(0,-4) -- (Z);
\draw<5-6>[celeste,<-,ultra thick,dashed](Y) ++(0,-4) -- (Y);


\fill(W) ++(4,0) circle(.1) coordinate(W);
\fill[rojo](X) ++(4,0) circle(.1) coordinate(X);
\fill(Y) ++(4,0) circle(.1) coordinate(Y);
\fill(Z) ++(4,0) circle(.1) coordinate(Z);
\fill(W) ++(0,-4) circle(.1) coordinate(W2);
\fill(Y) ++(0,-4) circle(.1) coordinate(Y2);
\fill(X) ++(0,4) circle(.1) coordinate(X2);
\fill(Z) ++(0,4) circle(.1) coordinate(Z2);
\draw[dashed](W2) circle(1);
\draw[dashed](X2) circle(1);
\draw[dashed](Y2) circle(1);
\draw[dashed](Z2) circle(1);
\draw[dashed](W) circle(1);
\draw[ultra thick, rojo,dashed](X) circle(1);
\draw[dashed](Y) circle(1);
\draw[dashed](Z) circle(1);

\fill(W) ++(4,0) circle(.1) coordinate(W);
\fill(X) ++(4,0) circle(.1) coordinate(X);
\fill(Y) ++(4,0) circle(.1) coordinate(Y);
\fill(Z) ++(4,0) circle(.1) coordinate(Z);
\fill(W) ++(0,-4) circle(.1) coordinate(W2);
\fill(Y) ++(0,-4) circle(.1) coordinate(Y2);
\fill(X) ++(0,4) circle(.1) coordinate(X2);
\fill(Z) ++(0,4) circle(.1) coordinate(Z2);
\draw[dashed](W2) circle(1);
\draw[dashed](X2) circle(1);
\draw[dashed](Y2) circle(1);
\draw[dashed](Z2) circle(1);
\draw[dashed](W) circle(1);
\draw[dashed](X) circle(1);
\draw[dashed](Y) circle(1);
\draw[dashed](Z) circle(1);

\fill(W) ++(-12,0) circle(.1) coordinate(W);
\fill(X) ++(-12,0) circle(.1) coordinate(X);
\fill(Y) ++(-12,0) circle(.1) coordinate(Y);
\fill[naranja](Z) ++(-12,0) circle(.1) coordinate(Z);
\fill(W) ++(0,-4) circle(.1) coordinate(W2);
\fill(Y) ++(0,-4) circle(.1) coordinate(Y2);
\fill(X) ++(0,4) circle(.1) coordinate(X2);
\fill(Z) ++(0,4) circle(.1) coordinate(Z2);
\draw[dashed](W2) circle(1);
\draw[dashed](X2) circle(1);
\draw[dashed](Y2) circle(1);
\draw[dashed](Z2) circle(1);
\draw[dashed](W) circle(1);
\draw[dashed](X) circle(1);
\draw[dashed](Y) circle(1);
\draw[ultra thick, naranja,dashed](Z) circle(1);
           
\end{scope} 
      

\end{zframe}
         
           
 
\end{document}
